\documentclass[11pt,a4paper]{article}

% shared/macros.tex
% Shared notation for all surface-partition math documents.
% Include with: % shared/macros.tex
% Shared notation for all surface-partition math documents.
% Include with: % shared/macros.tex
% Shared notation for all surface-partition math documents.
% Include with: \input{../shared/macros}

% -------------------------------------------------------------------------
% Packages (include here so each document only needs \input{../shared/macros})
% -------------------------------------------------------------------------
\usepackage{amsmath,amssymb,amsthm}
\usepackage{bm}           % \bm for bold math
\usepackage{mathtools}    % \coloneqq, dcases, etc.
\usepackage{booktabs}     % \toprule etc. in tables
\usepackage{hyperref}
\usepackage{geometry}
\usepackage{microtype}
\usepackage{parskip}
\usepackage{xcolor}
\usepackage{tcolorbox}
\usepackage{enumitem}
\usepackage{float}             % [H] float placement

\geometry{a4paper, margin=2.5cm}

% -------------------------------------------------------------------------
% Theorem-like environments
% -------------------------------------------------------------------------
\theoremstyle{definition}
\newtheorem{defn}{Definition}[section]
\newtheorem{remark}{Remark}[section]

\theoremstyle{plain}
\newtheorem{prop}{Proposition}[section]

% -------------------------------------------------------------------------
% Mesh and geometry
% -------------------------------------------------------------------------
\newcommand{\V}{\mathbf{V}}            % vertex array V[i] in R^dim
\newcommand{\F}{\mathbf{F}}            % face (triangle) array
\newcommand{\vone}[1]{v_{1,#1}}       % first vertex index of VP i
\newcommand{\vtwo}[1]{v_{2,#1}}       % second vertex index of VP i

% -------------------------------------------------------------------------
% Variable points
% -------------------------------------------------------------------------
\newcommand{\lam}{\lambda}             % scalar VP parameter
\newcommand{\Lam}{\boldsymbol{\lambda}}% full parameter vector
\newcommand{\vp}[1]{\mathbf{p}_{#1}}  % 3D position of VP i
\newcommand{\ed}[1]{\mathbf{d}_{#1}}  % edge direction d_i = V[v1_i] - V[v2_i]

% -------------------------------------------------------------------------
% Steiner / triple-point
% -------------------------------------------------------------------------
\newcommand{\Sv}{\mathbf{S}}           % Steiner (Fermat-Torricelli) point
\newcommand{\Ki}{K_{i}}               % projection matrix (I-n_i n_i^T)/r_i
\newcommand{\Mmat}{M}                 % sum of K_i matrices at a triple point

% -------------------------------------------------------------------------
% Normals and unit vectors
% -------------------------------------------------------------------------
\newcommand{\nhat}{\hat{\mathbf{n}}}  % unit normal to a triangle
\newcommand{\Dhat}{\hat{\boldsymbol{\Delta}}}  % unit segment direction

% -------------------------------------------------------------------------
% Operations
% -------------------------------------------------------------------------
\newcommand{\norm}[1]{\left\|#1\right\|}
\newcommand{\abs}[1]{\left|#1\right|}
\newcommand{\ip}[2]{\langle #1,\, #2 \rangle}

% Projection perpendicular to a unit vector \uhat
\newcommand{\Proj}[1]{\bigl(\mathbf{I} - #1 #1^\top\bigr)}

% argmax operator (winner-take-all assignment in Phase 1 documents)
\DeclareMathOperator*{\argmax}{arg\,max}
\DeclareMathOperator*{\argmin}{arg\,min}

% -------------------------------------------------------------------------
% Calculus shorthands
% -------------------------------------------------------------------------
\newcommand{\pd}[2]{\dfrac{\partial #1}{\partial #2}}
\newcommand{\pdd}[3]{\dfrac{\partial^2 #1}{\partial #2\,\partial #3}}
\newcommand{\grad}{\nabla}
\newcommand{\Hess}[1]{\nabla^2 #1}

% -------------------------------------------------------------------------
% Optimization problem objects
% -------------------------------------------------------------------------
\newcommand{\Preg}{P_{\mathrm{reg}}}  % regular perimeter
\newcommand{\Pst}{P_{\mathrm{st}}}   % Steiner perimeter
\newcommand{\Areg}{A^{\mathrm{reg}}} % regular cell area
\newcommand{\Ast}{A^{\mathrm{st}}}   % Steiner area contribution
\newcommand{\Atgt}{\bar{A}}           % target area per cell
\newcommand{\Lagr}{\mathcal{L}}       % Lagrangian
\newcommand{\objfac}{\sigma}          % IPOPT objective scaling factor

% -------------------------------------------------------------------------
% Code cross-references (for "In the code..." callouts)
% -------------------------------------------------------------------------
\newcommand{\code}[1]{\texttt{#1}}
\newcommand{\codefile}[1]{\texttt{\small #1}}

% -------------------------------------------------------------------------
% Threshold dynamics / approach B (10-mbo-auction-dynamics)
% -------------------------------------------------------------------------
\newcommand{\Dlump}{D}                 % lumped mass matrix diag(v)
\newcommand{\Atau}{A_\tau}             % discrete diffusion operator (M + tau K)^{-1} M
\newcommand{\Etau}{E_\tau}             % heat-content (threshold-dynamics) energy
\newcommand{\Rcell}{R_{\mathrm{cell}}} % radius of the equal-area disc sqrt(Abar/pi)
\newcommand{\hmean}{h_{\mathrm{mean}}} % mean edge length
\newcommand{\hmax}{h_{\mathrm{max}}}   % max edge length
\newcommand{\Gt}{G_t}                  % heat kernel at time t
\newcommand{\Ktau}{K^{\mathrm{res}}_\tau} % resolvent (backward-Euler) kernel


% -------------------------------------------------------------------------
% Packages (include here so each document only needs % shared/macros.tex
% Shared notation for all surface-partition math documents.
% Include with: \input{../shared/macros}

% -------------------------------------------------------------------------
% Packages (include here so each document only needs \input{../shared/macros})
% -------------------------------------------------------------------------
\usepackage{amsmath,amssymb,amsthm}
\usepackage{bm}           % \bm for bold math
\usepackage{mathtools}    % \coloneqq, dcases, etc.
\usepackage{booktabs}     % \toprule etc. in tables
\usepackage{hyperref}
\usepackage{geometry}
\usepackage{microtype}
\usepackage{parskip}
\usepackage{xcolor}
\usepackage{tcolorbox}
\usepackage{enumitem}
\usepackage{float}             % [H] float placement

\geometry{a4paper, margin=2.5cm}

% -------------------------------------------------------------------------
% Theorem-like environments
% -------------------------------------------------------------------------
\theoremstyle{definition}
\newtheorem{defn}{Definition}[section]
\newtheorem{remark}{Remark}[section]

\theoremstyle{plain}
\newtheorem{prop}{Proposition}[section]

% -------------------------------------------------------------------------
% Mesh and geometry
% -------------------------------------------------------------------------
\newcommand{\V}{\mathbf{V}}            % vertex array V[i] in R^dim
\newcommand{\F}{\mathbf{F}}            % face (triangle) array
\newcommand{\vone}[1]{v_{1,#1}}       % first vertex index of VP i
\newcommand{\vtwo}[1]{v_{2,#1}}       % second vertex index of VP i

% -------------------------------------------------------------------------
% Variable points
% -------------------------------------------------------------------------
\newcommand{\lam}{\lambda}             % scalar VP parameter
\newcommand{\Lam}{\boldsymbol{\lambda}}% full parameter vector
\newcommand{\vp}[1]{\mathbf{p}_{#1}}  % 3D position of VP i
\newcommand{\ed}[1]{\mathbf{d}_{#1}}  % edge direction d_i = V[v1_i] - V[v2_i]

% -------------------------------------------------------------------------
% Steiner / triple-point
% -------------------------------------------------------------------------
\newcommand{\Sv}{\mathbf{S}}           % Steiner (Fermat-Torricelli) point
\newcommand{\Ki}{K_{i}}               % projection matrix (I-n_i n_i^T)/r_i
\newcommand{\Mmat}{M}                 % sum of K_i matrices at a triple point

% -------------------------------------------------------------------------
% Normals and unit vectors
% -------------------------------------------------------------------------
\newcommand{\nhat}{\hat{\mathbf{n}}}  % unit normal to a triangle
\newcommand{\Dhat}{\hat{\boldsymbol{\Delta}}}  % unit segment direction

% -------------------------------------------------------------------------
% Operations
% -------------------------------------------------------------------------
\newcommand{\norm}[1]{\left\|#1\right\|}
\newcommand{\abs}[1]{\left|#1\right|}
\newcommand{\ip}[2]{\langle #1,\, #2 \rangle}

% Projection perpendicular to a unit vector \uhat
\newcommand{\Proj}[1]{\bigl(\mathbf{I} - #1 #1^\top\bigr)}

% argmax operator (winner-take-all assignment in Phase 1 documents)
\DeclareMathOperator*{\argmax}{arg\,max}
\DeclareMathOperator*{\argmin}{arg\,min}

% -------------------------------------------------------------------------
% Calculus shorthands
% -------------------------------------------------------------------------
\newcommand{\pd}[2]{\dfrac{\partial #1}{\partial #2}}
\newcommand{\pdd}[3]{\dfrac{\partial^2 #1}{\partial #2\,\partial #3}}
\newcommand{\grad}{\nabla}
\newcommand{\Hess}[1]{\nabla^2 #1}

% -------------------------------------------------------------------------
% Optimization problem objects
% -------------------------------------------------------------------------
\newcommand{\Preg}{P_{\mathrm{reg}}}  % regular perimeter
\newcommand{\Pst}{P_{\mathrm{st}}}   % Steiner perimeter
\newcommand{\Areg}{A^{\mathrm{reg}}} % regular cell area
\newcommand{\Ast}{A^{\mathrm{st}}}   % Steiner area contribution
\newcommand{\Atgt}{\bar{A}}           % target area per cell
\newcommand{\Lagr}{\mathcal{L}}       % Lagrangian
\newcommand{\objfac}{\sigma}          % IPOPT objective scaling factor

% -------------------------------------------------------------------------
% Code cross-references (for "In the code..." callouts)
% -------------------------------------------------------------------------
\newcommand{\code}[1]{\texttt{#1}}
\newcommand{\codefile}[1]{\texttt{\small #1}}

% -------------------------------------------------------------------------
% Threshold dynamics / approach B (10-mbo-auction-dynamics)
% -------------------------------------------------------------------------
\newcommand{\Dlump}{D}                 % lumped mass matrix diag(v)
\newcommand{\Atau}{A_\tau}             % discrete diffusion operator (M + tau K)^{-1} M
\newcommand{\Etau}{E_\tau}             % heat-content (threshold-dynamics) energy
\newcommand{\Rcell}{R_{\mathrm{cell}}} % radius of the equal-area disc sqrt(Abar/pi)
\newcommand{\hmean}{h_{\mathrm{mean}}} % mean edge length
\newcommand{\hmax}{h_{\mathrm{max}}}   % max edge length
\newcommand{\Gt}{G_t}                  % heat kernel at time t
\newcommand{\Ktau}{K^{\mathrm{res}}_\tau} % resolvent (backward-Euler) kernel
)
% -------------------------------------------------------------------------
\usepackage{amsmath,amssymb,amsthm}
\usepackage{bm}           % \bm for bold math
\usepackage{mathtools}    % \coloneqq, dcases, etc.
\usepackage{booktabs}     % \toprule etc. in tables
\usepackage{hyperref}
\usepackage{geometry}
\usepackage{microtype}
\usepackage{parskip}
\usepackage{xcolor}
\usepackage{tcolorbox}
\usepackage{enumitem}
\usepackage{float}             % [H] float placement

\geometry{a4paper, margin=2.5cm}

% -------------------------------------------------------------------------
% Theorem-like environments
% -------------------------------------------------------------------------
\theoremstyle{definition}
\newtheorem{defn}{Definition}[section]
\newtheorem{remark}{Remark}[section]

\theoremstyle{plain}
\newtheorem{prop}{Proposition}[section]

% -------------------------------------------------------------------------
% Mesh and geometry
% -------------------------------------------------------------------------
\newcommand{\V}{\mathbf{V}}            % vertex array V[i] in R^dim
\newcommand{\F}{\mathbf{F}}            % face (triangle) array
\newcommand{\vone}[1]{v_{1,#1}}       % first vertex index of VP i
\newcommand{\vtwo}[1]{v_{2,#1}}       % second vertex index of VP i

% -------------------------------------------------------------------------
% Variable points
% -------------------------------------------------------------------------
\newcommand{\lam}{\lambda}             % scalar VP parameter
\newcommand{\Lam}{\boldsymbol{\lambda}}% full parameter vector
\newcommand{\vp}[1]{\mathbf{p}_{#1}}  % 3D position of VP i
\newcommand{\ed}[1]{\mathbf{d}_{#1}}  % edge direction d_i = V[v1_i] - V[v2_i]

% -------------------------------------------------------------------------
% Steiner / triple-point
% -------------------------------------------------------------------------
\newcommand{\Sv}{\mathbf{S}}           % Steiner (Fermat-Torricelli) point
\newcommand{\Ki}{K_{i}}               % projection matrix (I-n_i n_i^T)/r_i
\newcommand{\Mmat}{M}                 % sum of K_i matrices at a triple point

% -------------------------------------------------------------------------
% Normals and unit vectors
% -------------------------------------------------------------------------
\newcommand{\nhat}{\hat{\mathbf{n}}}  % unit normal to a triangle
\newcommand{\Dhat}{\hat{\boldsymbol{\Delta}}}  % unit segment direction

% -------------------------------------------------------------------------
% Operations
% -------------------------------------------------------------------------
\newcommand{\norm}[1]{\left\|#1\right\|}
\newcommand{\abs}[1]{\left|#1\right|}
\newcommand{\ip}[2]{\langle #1,\, #2 \rangle}

% Projection perpendicular to a unit vector \uhat
\newcommand{\Proj}[1]{\bigl(\mathbf{I} - #1 #1^\top\bigr)}

% argmax operator (winner-take-all assignment in Phase 1 documents)
\DeclareMathOperator*{\argmax}{arg\,max}
\DeclareMathOperator*{\argmin}{arg\,min}

% -------------------------------------------------------------------------
% Calculus shorthands
% -------------------------------------------------------------------------
\newcommand{\pd}[2]{\dfrac{\partial #1}{\partial #2}}
\newcommand{\pdd}[3]{\dfrac{\partial^2 #1}{\partial #2\,\partial #3}}
\newcommand{\grad}{\nabla}
\newcommand{\Hess}[1]{\nabla^2 #1}

% -------------------------------------------------------------------------
% Optimization problem objects
% -------------------------------------------------------------------------
\newcommand{\Preg}{P_{\mathrm{reg}}}  % regular perimeter
\newcommand{\Pst}{P_{\mathrm{st}}}   % Steiner perimeter
\newcommand{\Areg}{A^{\mathrm{reg}}} % regular cell area
\newcommand{\Ast}{A^{\mathrm{st}}}   % Steiner area contribution
\newcommand{\Atgt}{\bar{A}}           % target area per cell
\newcommand{\Lagr}{\mathcal{L}}       % Lagrangian
\newcommand{\objfac}{\sigma}          % IPOPT objective scaling factor

% -------------------------------------------------------------------------
% Code cross-references (for "In the code..." callouts)
% -------------------------------------------------------------------------
\newcommand{\code}[1]{\texttt{#1}}
\newcommand{\codefile}[1]{\texttt{\small #1}}

% -------------------------------------------------------------------------
% Threshold dynamics / approach B (10-mbo-auction-dynamics)
% -------------------------------------------------------------------------
\newcommand{\Dlump}{D}                 % lumped mass matrix diag(v)
\newcommand{\Atau}{A_\tau}             % discrete diffusion operator (M + tau K)^{-1} M
\newcommand{\Etau}{E_\tau}             % heat-content (threshold-dynamics) energy
\newcommand{\Rcell}{R_{\mathrm{cell}}} % radius of the equal-area disc sqrt(Abar/pi)
\newcommand{\hmean}{h_{\mathrm{mean}}} % mean edge length
\newcommand{\hmax}{h_{\mathrm{max}}}   % max edge length
\newcommand{\Gt}{G_t}                  % heat kernel at time t
\newcommand{\Ktau}{K^{\mathrm{res}}_\tau} % resolvent (backward-Euler) kernel


% -------------------------------------------------------------------------
% Packages (include here so each document only needs % shared/macros.tex
% Shared notation for all surface-partition math documents.
% Include with: % shared/macros.tex
% Shared notation for all surface-partition math documents.
% Include with: \input{../shared/macros}

% -------------------------------------------------------------------------
% Packages (include here so each document only needs \input{../shared/macros})
% -------------------------------------------------------------------------
\usepackage{amsmath,amssymb,amsthm}
\usepackage{bm}           % \bm for bold math
\usepackage{mathtools}    % \coloneqq, dcases, etc.
\usepackage{booktabs}     % \toprule etc. in tables
\usepackage{hyperref}
\usepackage{geometry}
\usepackage{microtype}
\usepackage{parskip}
\usepackage{xcolor}
\usepackage{tcolorbox}
\usepackage{enumitem}
\usepackage{float}             % [H] float placement

\geometry{a4paper, margin=2.5cm}

% -------------------------------------------------------------------------
% Theorem-like environments
% -------------------------------------------------------------------------
\theoremstyle{definition}
\newtheorem{defn}{Definition}[section]
\newtheorem{remark}{Remark}[section]

\theoremstyle{plain}
\newtheorem{prop}{Proposition}[section]

% -------------------------------------------------------------------------
% Mesh and geometry
% -------------------------------------------------------------------------
\newcommand{\V}{\mathbf{V}}            % vertex array V[i] in R^dim
\newcommand{\F}{\mathbf{F}}            % face (triangle) array
\newcommand{\vone}[1]{v_{1,#1}}       % first vertex index of VP i
\newcommand{\vtwo}[1]{v_{2,#1}}       % second vertex index of VP i

% -------------------------------------------------------------------------
% Variable points
% -------------------------------------------------------------------------
\newcommand{\lam}{\lambda}             % scalar VP parameter
\newcommand{\Lam}{\boldsymbol{\lambda}}% full parameter vector
\newcommand{\vp}[1]{\mathbf{p}_{#1}}  % 3D position of VP i
\newcommand{\ed}[1]{\mathbf{d}_{#1}}  % edge direction d_i = V[v1_i] - V[v2_i]

% -------------------------------------------------------------------------
% Steiner / triple-point
% -------------------------------------------------------------------------
\newcommand{\Sv}{\mathbf{S}}           % Steiner (Fermat-Torricelli) point
\newcommand{\Ki}{K_{i}}               % projection matrix (I-n_i n_i^T)/r_i
\newcommand{\Mmat}{M}                 % sum of K_i matrices at a triple point

% -------------------------------------------------------------------------
% Normals and unit vectors
% -------------------------------------------------------------------------
\newcommand{\nhat}{\hat{\mathbf{n}}}  % unit normal to a triangle
\newcommand{\Dhat}{\hat{\boldsymbol{\Delta}}}  % unit segment direction

% -------------------------------------------------------------------------
% Operations
% -------------------------------------------------------------------------
\newcommand{\norm}[1]{\left\|#1\right\|}
\newcommand{\abs}[1]{\left|#1\right|}
\newcommand{\ip}[2]{\langle #1,\, #2 \rangle}

% Projection perpendicular to a unit vector \uhat
\newcommand{\Proj}[1]{\bigl(\mathbf{I} - #1 #1^\top\bigr)}

% argmax operator (winner-take-all assignment in Phase 1 documents)
\DeclareMathOperator*{\argmax}{arg\,max}
\DeclareMathOperator*{\argmin}{arg\,min}

% -------------------------------------------------------------------------
% Calculus shorthands
% -------------------------------------------------------------------------
\newcommand{\pd}[2]{\dfrac{\partial #1}{\partial #2}}
\newcommand{\pdd}[3]{\dfrac{\partial^2 #1}{\partial #2\,\partial #3}}
\newcommand{\grad}{\nabla}
\newcommand{\Hess}[1]{\nabla^2 #1}

% -------------------------------------------------------------------------
% Optimization problem objects
% -------------------------------------------------------------------------
\newcommand{\Preg}{P_{\mathrm{reg}}}  % regular perimeter
\newcommand{\Pst}{P_{\mathrm{st}}}   % Steiner perimeter
\newcommand{\Areg}{A^{\mathrm{reg}}} % regular cell area
\newcommand{\Ast}{A^{\mathrm{st}}}   % Steiner area contribution
\newcommand{\Atgt}{\bar{A}}           % target area per cell
\newcommand{\Lagr}{\mathcal{L}}       % Lagrangian
\newcommand{\objfac}{\sigma}          % IPOPT objective scaling factor

% -------------------------------------------------------------------------
% Code cross-references (for "In the code..." callouts)
% -------------------------------------------------------------------------
\newcommand{\code}[1]{\texttt{#1}}
\newcommand{\codefile}[1]{\texttt{\small #1}}

% -------------------------------------------------------------------------
% Threshold dynamics / approach B (10-mbo-auction-dynamics)
% -------------------------------------------------------------------------
\newcommand{\Dlump}{D}                 % lumped mass matrix diag(v)
\newcommand{\Atau}{A_\tau}             % discrete diffusion operator (M + tau K)^{-1} M
\newcommand{\Etau}{E_\tau}             % heat-content (threshold-dynamics) energy
\newcommand{\Rcell}{R_{\mathrm{cell}}} % radius of the equal-area disc sqrt(Abar/pi)
\newcommand{\hmean}{h_{\mathrm{mean}}} % mean edge length
\newcommand{\hmax}{h_{\mathrm{max}}}   % max edge length
\newcommand{\Gt}{G_t}                  % heat kernel at time t
\newcommand{\Ktau}{K^{\mathrm{res}}_\tau} % resolvent (backward-Euler) kernel


% -------------------------------------------------------------------------
% Packages (include here so each document only needs % shared/macros.tex
% Shared notation for all surface-partition math documents.
% Include with: \input{../shared/macros}

% -------------------------------------------------------------------------
% Packages (include here so each document only needs \input{../shared/macros})
% -------------------------------------------------------------------------
\usepackage{amsmath,amssymb,amsthm}
\usepackage{bm}           % \bm for bold math
\usepackage{mathtools}    % \coloneqq, dcases, etc.
\usepackage{booktabs}     % \toprule etc. in tables
\usepackage{hyperref}
\usepackage{geometry}
\usepackage{microtype}
\usepackage{parskip}
\usepackage{xcolor}
\usepackage{tcolorbox}
\usepackage{enumitem}
\usepackage{float}             % [H] float placement

\geometry{a4paper, margin=2.5cm}

% -------------------------------------------------------------------------
% Theorem-like environments
% -------------------------------------------------------------------------
\theoremstyle{definition}
\newtheorem{defn}{Definition}[section]
\newtheorem{remark}{Remark}[section]

\theoremstyle{plain}
\newtheorem{prop}{Proposition}[section]

% -------------------------------------------------------------------------
% Mesh and geometry
% -------------------------------------------------------------------------
\newcommand{\V}{\mathbf{V}}            % vertex array V[i] in R^dim
\newcommand{\F}{\mathbf{F}}            % face (triangle) array
\newcommand{\vone}[1]{v_{1,#1}}       % first vertex index of VP i
\newcommand{\vtwo}[1]{v_{2,#1}}       % second vertex index of VP i

% -------------------------------------------------------------------------
% Variable points
% -------------------------------------------------------------------------
\newcommand{\lam}{\lambda}             % scalar VP parameter
\newcommand{\Lam}{\boldsymbol{\lambda}}% full parameter vector
\newcommand{\vp}[1]{\mathbf{p}_{#1}}  % 3D position of VP i
\newcommand{\ed}[1]{\mathbf{d}_{#1}}  % edge direction d_i = V[v1_i] - V[v2_i]

% -------------------------------------------------------------------------
% Steiner / triple-point
% -------------------------------------------------------------------------
\newcommand{\Sv}{\mathbf{S}}           % Steiner (Fermat-Torricelli) point
\newcommand{\Ki}{K_{i}}               % projection matrix (I-n_i n_i^T)/r_i
\newcommand{\Mmat}{M}                 % sum of K_i matrices at a triple point

% -------------------------------------------------------------------------
% Normals and unit vectors
% -------------------------------------------------------------------------
\newcommand{\nhat}{\hat{\mathbf{n}}}  % unit normal to a triangle
\newcommand{\Dhat}{\hat{\boldsymbol{\Delta}}}  % unit segment direction

% -------------------------------------------------------------------------
% Operations
% -------------------------------------------------------------------------
\newcommand{\norm}[1]{\left\|#1\right\|}
\newcommand{\abs}[1]{\left|#1\right|}
\newcommand{\ip}[2]{\langle #1,\, #2 \rangle}

% Projection perpendicular to a unit vector \uhat
\newcommand{\Proj}[1]{\bigl(\mathbf{I} - #1 #1^\top\bigr)}

% argmax operator (winner-take-all assignment in Phase 1 documents)
\DeclareMathOperator*{\argmax}{arg\,max}
\DeclareMathOperator*{\argmin}{arg\,min}

% -------------------------------------------------------------------------
% Calculus shorthands
% -------------------------------------------------------------------------
\newcommand{\pd}[2]{\dfrac{\partial #1}{\partial #2}}
\newcommand{\pdd}[3]{\dfrac{\partial^2 #1}{\partial #2\,\partial #3}}
\newcommand{\grad}{\nabla}
\newcommand{\Hess}[1]{\nabla^2 #1}

% -------------------------------------------------------------------------
% Optimization problem objects
% -------------------------------------------------------------------------
\newcommand{\Preg}{P_{\mathrm{reg}}}  % regular perimeter
\newcommand{\Pst}{P_{\mathrm{st}}}   % Steiner perimeter
\newcommand{\Areg}{A^{\mathrm{reg}}} % regular cell area
\newcommand{\Ast}{A^{\mathrm{st}}}   % Steiner area contribution
\newcommand{\Atgt}{\bar{A}}           % target area per cell
\newcommand{\Lagr}{\mathcal{L}}       % Lagrangian
\newcommand{\objfac}{\sigma}          % IPOPT objective scaling factor

% -------------------------------------------------------------------------
% Code cross-references (for "In the code..." callouts)
% -------------------------------------------------------------------------
\newcommand{\code}[1]{\texttt{#1}}
\newcommand{\codefile}[1]{\texttt{\small #1}}

% -------------------------------------------------------------------------
% Threshold dynamics / approach B (10-mbo-auction-dynamics)
% -------------------------------------------------------------------------
\newcommand{\Dlump}{D}                 % lumped mass matrix diag(v)
\newcommand{\Atau}{A_\tau}             % discrete diffusion operator (M + tau K)^{-1} M
\newcommand{\Etau}{E_\tau}             % heat-content (threshold-dynamics) energy
\newcommand{\Rcell}{R_{\mathrm{cell}}} % radius of the equal-area disc sqrt(Abar/pi)
\newcommand{\hmean}{h_{\mathrm{mean}}} % mean edge length
\newcommand{\hmax}{h_{\mathrm{max}}}   % max edge length
\newcommand{\Gt}{G_t}                  % heat kernel at time t
\newcommand{\Ktau}{K^{\mathrm{res}}_\tau} % resolvent (backward-Euler) kernel
)
% -------------------------------------------------------------------------
\usepackage{amsmath,amssymb,amsthm}
\usepackage{bm}           % \bm for bold math
\usepackage{mathtools}    % \coloneqq, dcases, etc.
\usepackage{booktabs}     % \toprule etc. in tables
\usepackage{hyperref}
\usepackage{geometry}
\usepackage{microtype}
\usepackage{parskip}
\usepackage{xcolor}
\usepackage{tcolorbox}
\usepackage{enumitem}
\usepackage{float}             % [H] float placement

\geometry{a4paper, margin=2.5cm}

% -------------------------------------------------------------------------
% Theorem-like environments
% -------------------------------------------------------------------------
\theoremstyle{definition}
\newtheorem{defn}{Definition}[section]
\newtheorem{remark}{Remark}[section]

\theoremstyle{plain}
\newtheorem{prop}{Proposition}[section]

% -------------------------------------------------------------------------
% Mesh and geometry
% -------------------------------------------------------------------------
\newcommand{\V}{\mathbf{V}}            % vertex array V[i] in R^dim
\newcommand{\F}{\mathbf{F}}            % face (triangle) array
\newcommand{\vone}[1]{v_{1,#1}}       % first vertex index of VP i
\newcommand{\vtwo}[1]{v_{2,#1}}       % second vertex index of VP i

% -------------------------------------------------------------------------
% Variable points
% -------------------------------------------------------------------------
\newcommand{\lam}{\lambda}             % scalar VP parameter
\newcommand{\Lam}{\boldsymbol{\lambda}}% full parameter vector
\newcommand{\vp}[1]{\mathbf{p}_{#1}}  % 3D position of VP i
\newcommand{\ed}[1]{\mathbf{d}_{#1}}  % edge direction d_i = V[v1_i] - V[v2_i]

% -------------------------------------------------------------------------
% Steiner / triple-point
% -------------------------------------------------------------------------
\newcommand{\Sv}{\mathbf{S}}           % Steiner (Fermat-Torricelli) point
\newcommand{\Ki}{K_{i}}               % projection matrix (I-n_i n_i^T)/r_i
\newcommand{\Mmat}{M}                 % sum of K_i matrices at a triple point

% -------------------------------------------------------------------------
% Normals and unit vectors
% -------------------------------------------------------------------------
\newcommand{\nhat}{\hat{\mathbf{n}}}  % unit normal to a triangle
\newcommand{\Dhat}{\hat{\boldsymbol{\Delta}}}  % unit segment direction

% -------------------------------------------------------------------------
% Operations
% -------------------------------------------------------------------------
\newcommand{\norm}[1]{\left\|#1\right\|}
\newcommand{\abs}[1]{\left|#1\right|}
\newcommand{\ip}[2]{\langle #1,\, #2 \rangle}

% Projection perpendicular to a unit vector \uhat
\newcommand{\Proj}[1]{\bigl(\mathbf{I} - #1 #1^\top\bigr)}

% argmax operator (winner-take-all assignment in Phase 1 documents)
\DeclareMathOperator*{\argmax}{arg\,max}
\DeclareMathOperator*{\argmin}{arg\,min}

% -------------------------------------------------------------------------
% Calculus shorthands
% -------------------------------------------------------------------------
\newcommand{\pd}[2]{\dfrac{\partial #1}{\partial #2}}
\newcommand{\pdd}[3]{\dfrac{\partial^2 #1}{\partial #2\,\partial #3}}
\newcommand{\grad}{\nabla}
\newcommand{\Hess}[1]{\nabla^2 #1}

% -------------------------------------------------------------------------
% Optimization problem objects
% -------------------------------------------------------------------------
\newcommand{\Preg}{P_{\mathrm{reg}}}  % regular perimeter
\newcommand{\Pst}{P_{\mathrm{st}}}   % Steiner perimeter
\newcommand{\Areg}{A^{\mathrm{reg}}} % regular cell area
\newcommand{\Ast}{A^{\mathrm{st}}}   % Steiner area contribution
\newcommand{\Atgt}{\bar{A}}           % target area per cell
\newcommand{\Lagr}{\mathcal{L}}       % Lagrangian
\newcommand{\objfac}{\sigma}          % IPOPT objective scaling factor

% -------------------------------------------------------------------------
% Code cross-references (for "In the code..." callouts)
% -------------------------------------------------------------------------
\newcommand{\code}[1]{\texttt{#1}}
\newcommand{\codefile}[1]{\texttt{\small #1}}

% -------------------------------------------------------------------------
% Threshold dynamics / approach B (10-mbo-auction-dynamics)
% -------------------------------------------------------------------------
\newcommand{\Dlump}{D}                 % lumped mass matrix diag(v)
\newcommand{\Atau}{A_\tau}             % discrete diffusion operator (M + tau K)^{-1} M
\newcommand{\Etau}{E_\tau}             % heat-content (threshold-dynamics) energy
\newcommand{\Rcell}{R_{\mathrm{cell}}} % radius of the equal-area disc sqrt(Abar/pi)
\newcommand{\hmean}{h_{\mathrm{mean}}} % mean edge length
\newcommand{\hmax}{h_{\mathrm{max}}}   % max edge length
\newcommand{\Gt}{G_t}                  % heat kernel at time t
\newcommand{\Ktau}{K^{\mathrm{res}}_\tau} % resolvent (backward-Euler) kernel
)
% -------------------------------------------------------------------------
\usepackage{amsmath,amssymb,amsthm}
\usepackage{bm}           % \bm for bold math
\usepackage{mathtools}    % \coloneqq, dcases, etc.
\usepackage{booktabs}     % \toprule etc. in tables
\usepackage{hyperref}
\usepackage{geometry}
\usepackage{microtype}
\usepackage{parskip}
\usepackage{xcolor}
\usepackage{tcolorbox}
\usepackage{enumitem}
\usepackage{float}             % [H] float placement

\geometry{a4paper, margin=2.5cm}

% -------------------------------------------------------------------------
% Theorem-like environments
% -------------------------------------------------------------------------
\theoremstyle{definition}
\newtheorem{defn}{Definition}[section]
\newtheorem{remark}{Remark}[section]

\theoremstyle{plain}
\newtheorem{prop}{Proposition}[section]

% -------------------------------------------------------------------------
% Mesh and geometry
% -------------------------------------------------------------------------
\newcommand{\V}{\mathbf{V}}            % vertex array V[i] in R^dim
\newcommand{\F}{\mathbf{F}}            % face (triangle) array
\newcommand{\vone}[1]{v_{1,#1}}       % first vertex index of VP i
\newcommand{\vtwo}[1]{v_{2,#1}}       % second vertex index of VP i

% -------------------------------------------------------------------------
% Variable points
% -------------------------------------------------------------------------
\newcommand{\lam}{\lambda}             % scalar VP parameter
\newcommand{\Lam}{\boldsymbol{\lambda}}% full parameter vector
\newcommand{\vp}[1]{\mathbf{p}_{#1}}  % 3D position of VP i
\newcommand{\ed}[1]{\mathbf{d}_{#1}}  % edge direction d_i = V[v1_i] - V[v2_i]

% -------------------------------------------------------------------------
% Steiner / triple-point
% -------------------------------------------------------------------------
\newcommand{\Sv}{\mathbf{S}}           % Steiner (Fermat-Torricelli) point
\newcommand{\Ki}{K_{i}}               % projection matrix (I-n_i n_i^T)/r_i
\newcommand{\Mmat}{M}                 % sum of K_i matrices at a triple point

% -------------------------------------------------------------------------
% Normals and unit vectors
% -------------------------------------------------------------------------
\newcommand{\nhat}{\hat{\mathbf{n}}}  % unit normal to a triangle
\newcommand{\Dhat}{\hat{\boldsymbol{\Delta}}}  % unit segment direction

% -------------------------------------------------------------------------
% Operations
% -------------------------------------------------------------------------
\newcommand{\norm}[1]{\left\|#1\right\|}
\newcommand{\abs}[1]{\left|#1\right|}
\newcommand{\ip}[2]{\langle #1,\, #2 \rangle}

% Projection perpendicular to a unit vector \uhat
\newcommand{\Proj}[1]{\bigl(\mathbf{I} - #1 #1^\top\bigr)}

% argmax operator (winner-take-all assignment in Phase 1 documents)
\DeclareMathOperator*{\argmax}{arg\,max}
\DeclareMathOperator*{\argmin}{arg\,min}

% -------------------------------------------------------------------------
% Calculus shorthands
% -------------------------------------------------------------------------
\newcommand{\pd}[2]{\dfrac{\partial #1}{\partial #2}}
\newcommand{\pdd}[3]{\dfrac{\partial^2 #1}{\partial #2\,\partial #3}}
\newcommand{\grad}{\nabla}
\newcommand{\Hess}[1]{\nabla^2 #1}

% -------------------------------------------------------------------------
% Optimization problem objects
% -------------------------------------------------------------------------
\newcommand{\Preg}{P_{\mathrm{reg}}}  % regular perimeter
\newcommand{\Pst}{P_{\mathrm{st}}}   % Steiner perimeter
\newcommand{\Areg}{A^{\mathrm{reg}}} % regular cell area
\newcommand{\Ast}{A^{\mathrm{st}}}   % Steiner area contribution
\newcommand{\Atgt}{\bar{A}}           % target area per cell
\newcommand{\Lagr}{\mathcal{L}}       % Lagrangian
\newcommand{\objfac}{\sigma}          % IPOPT objective scaling factor

% -------------------------------------------------------------------------
% Code cross-references (for "In the code..." callouts)
% -------------------------------------------------------------------------
\newcommand{\code}[1]{\texttt{#1}}
\newcommand{\codefile}[1]{\texttt{\small #1}}

% -------------------------------------------------------------------------
% Threshold dynamics / approach B (10-mbo-auction-dynamics)
% -------------------------------------------------------------------------
\newcommand{\Dlump}{D}                 % lumped mass matrix diag(v)
\newcommand{\Atau}{A_\tau}             % discrete diffusion operator (M + tau K)^{-1} M
\newcommand{\Etau}{E_\tau}             % heat-content (threshold-dynamics) energy
\newcommand{\Rcell}{R_{\mathrm{cell}}} % radius of the equal-area disc sqrt(Abar/pi)
\newcommand{\hmean}{h_{\mathrm{mean}}} % mean edge length
\newcommand{\hmax}{h_{\mathrm{max}}}   % max edge length
\newcommand{\Gt}{G_t}                  % heat kernel at time t
\newcommand{\Ktau}{K^{\mathrm{res}}_\tau} % resolvent (backward-Euler) kernel
      % loads ALL packages and notation

\hypersetup{
  colorlinks = true,
  linkcolor  = blue!60!black,
  citecolor  = green!50!black,
  urlcolor   = blue!60!black,
  pdftitle   = {Analytical Steiner Derivatives},
  pdfauthor  = {surface-partition project}
}

\title{\textbf{Analytical Steiner Derivatives}\\[0.5em]
  \large Closed-form first and second derivatives of triple-point contributions}
\author{}
\date{May 2026}

\begin{document}
\maketitle

\begin{abstract}
At a triple point, three partition boundaries meet at a Fermat--Torricelli
(Steiner) point $\Sv$.  The Phase~2 perimeter-refinement optimizer needs the
first and second derivatives of the Steiner perimeter and area contributions
with respect to the variable-point parameters $\Lam$.  This document derives
all of them in closed form by applying the implicit-function theorem to the
Fermat-point optimality condition.  Every quantity here is implemented
analytically in \codefile{src/partition/vectorized\_steiner.py}; the companion
document \codefile{docs/math/01-phase2-derivatives} covers the regular
(non-Steiner) perimeter and area derivatives, and the finite-difference
schemes that these closed forms replaced.
\end{abstract}

\tableofcontents
\newpage

% =========================================================================
\section{Overview}
% =========================================================================

A \emph{triple point} arises at a mesh triangle where three partition cells
meet.  Inside that triangle the three cell boundaries do not run straight to a
shared mesh vertex; instead they meet at the \emph{Fermat--Torricelli point}
$\Sv$ --- the point minimising the total distance to the triangle's three
variable points~\cite{kuhn1974steiner}.  The Steiner construction contributes
to both the perimeter objective and the per-cell area constraints, so the
optimizer needs

\begin{equation}
  \pd{\Pst}{\lam_i},\qquad
  \pd{\Ast_k}{\lam_i},\qquad
  \pdd{\Pst}{\lam_i}{\lam_j},\qquad
  \pdd{\Ast_k}{\lam_i}{\lam_j}.
\end{equation}

The forward value $\Sv$ is a closed-form barycentric expression, but its
derivatives are not elementary: $\Sv$ is defined only implicitly, as the
solution of an optimality condition.  The key is the implicit-function
theorem.  Once $\partial\Sv/\partial\lam_i$ and
$\partial^2\Sv/\partial\lam_i\partial\lam_j$ are known, the perimeter and area
derivatives follow by the chain rule.

Sections~\ref{sec:setup}--\ref{sec:first-S} derive $\partial\Sv/\partial p_k$;
Section~\ref{sec:first} gives the first derivatives of the perimeter and area;
Sections~\ref{sec:second-S}--\ref{sec:lambda} derive
$\partial^2\Sv/\partial p_k\partial p_l$ and lift it to $\lam$-space;
Section~\ref{sec:hessians} assembles the perimeter and area Hessians;
Section~\ref{sec:degenerate} treats the degenerate case.

% =========================================================================
\section{Notation and Setup}
\label{sec:setup}
% =========================================================================

At one triple point three variable points $\vp{1},\vp{2},\vp{3}\in\mathbb{R}^3$
each depend affinely on a scalar parameter:
\begin{equation}
  \vp{i}(\lam_i) = \lam_i\,\V[\vone{i}] + (1-\lam_i)\,\V[\vtwo{i}],
  \qquad
  \ed{i} \coloneqq \pd{\vp{i}}{\lam_i} = \V[\vone{i}] - \V[\vtwo{i}].
\end{equation}
The edge direction $\ed{i}$ is constant and $\partial^2\vp{i}/\partial\lam_i^2
= \mathbf{0}$: every variable point is \emph{linear} in its own parameter.
This is the single structural fact that makes the chain rule below tractable.

Let $\Sv$ be the Fermat--Torricelli point of the triangle
$(\vp{1},\vp{2},\vp{3})$.  Define, for $i = 1,2,3$,
\begin{equation}
  \mathbf{u}_i \coloneqq \vp{i} - \Sv,\qquad
  r_i \coloneqq \norm{\mathbf{u}_i},\qquad
  \mathbf{n}_i \coloneqq \mathbf{u}_i / r_i,
\end{equation}
\begin{equation}
  K_i \coloneqq \frac{\Proj{\mathbf{n}_i}}{r_i},\qquad
  \Mmat \coloneqq K_1 + K_2 + K_3.
  \label{eq:KM}
\end{equation}
Each $K_i$ is symmetric positive semidefinite; $\Mmat$ is symmetric positive
definite whenever the triangle is non-degenerate (Section~\ref{sec:degenerate}).

\begin{tcolorbox}[colback=blue!4!white, colframe=blue!40!black,
                  title={Code: triple-point geometry}, fontupper=\small]
\codefile{src/partition/vectorized\_steiner.py}:
\code{\_compute\_tp\_geometry(pa)} returns $(\vp{}, \Sv, r, \mathbf{n}, K,
\Mmat)$ per triple point; \code{\_compute\_tp\_steiner(pa)} is the shared
forward-value / degeneracy core.
\end{tcolorbox}

% =========================================================================
\section{The Fermat Point as an Implicit Equation}
\label{sec:implicit}
% =========================================================================

$\Sv$ is the unique minimiser of the strictly convex function
\begin{equation}
  G(\Sv) = \norm{\vp{1}-\Sv} + \norm{\vp{2}-\Sv} + \norm{\vp{3}-\Sv}.
\end{equation}
Since $\partial\norm{\vp{i}-\Sv}/\partial\Sv = -\mathbf{n}_i$, the first-order
optimality condition $\grad_{\Sv} G = \mathbf{0}$ is
\begin{equation}
  F(\Sv,\vp{1},\vp{2},\vp{3}) \;\coloneqq\; \sum_{i=1}^{3}\mathbf{n}_i
  \;=\; \mathbf{0}.
  \label{eq:fermat}
\end{equation}
Geometrically the three unit vectors from $\Sv$ to the variable points sum to
zero --- equivalently they make $120^\circ$ angles, the defining property of
the Fermat point when every triangle angle is below $120^\circ$.  Equation
\eqref{eq:fermat} defines $\Sv$ implicitly; differentiating it is the route to
all the derivatives below.

\begin{prop}[Derivative of a unit vector]
\label{prop:unit}
If $\mathbf{u}$ depends on a variable $\mathbf{x}$ with constant Jacobian
$\partial\mathbf{u}/\partial\mathbf{x} = A$, then the unit vector
$\mathbf{n}=\mathbf{u}/r$, $r=\norm{\mathbf{u}}$, satisfies
\begin{equation}
  \pd{\mathbf{n}}{\mathbf{x}}
   = \frac{\Proj{\mathbf{n}}}{r}\,A
   = K\,A,
  \qquad
  \pd{r}{\mathbf{x}} = \mathbf{n}^\top A .
  \label{eq:unit-deriv}
\end{equation}
\end{prop}
\begin{proof}
$\partial r/\partial\mathbf{x} = (\mathbf{u}^\top/r)\,A = \mathbf{n}^\top A$.
Then $\partial\mathbf{n}/\partial\mathbf{x}
= A/r - (\mathbf{u}/r^2)\,\partial r/\partial\mathbf{x}
= \bigl(\mathbf{I}-\mathbf{n}\mathbf{n}^\top\bigr)A/r = K A$.
\end{proof}

% =========================================================================
\section{First Derivative of the Steiner Point}
\label{sec:first-S}
% =========================================================================

Differentiate the optimality condition \eqref{eq:fermat} with respect to
$\vp{k}$.  By the implicit-function theorem,
\begin{equation}
  \pd{F}{\Sv}\,\pd{\Sv}{\vp{k}} + \pd{F}{\vp{k}} = \mathbf{0}.
\end{equation}
Using Proposition~\ref{prop:unit} with $\mathbf{u}_i = \vp{i}-\Sv$:
varying $\Sv$ gives $\partial\mathbf{u}_i/\partial\Sv = -\mathbf{I}$, so
$\partial F/\partial\Sv = \sum_i K_i(-\mathbf{I}) = -\Mmat$; varying $\vp{k}$
gives $\partial\mathbf{u}_i/\partial\vp{k} = \delta_{ik}\mathbf{I}$, so
$\partial F/\partial\vp{k} = K_k$.  Hence
\begin{equation}
  -\Mmat\,\pd{\Sv}{\vp{k}} + K_k = \mathbf{0}
  \quad\Longrightarrow\quad
  \boxed{\;\pd{\Sv}{\vp{k}} = \Mmat^{-1}K_k\;}
  \label{eq:dS-dp}
\end{equation}

\begin{remark}[Translation invariance]
\label{rem:translation}
Summing \eqref{eq:dS-dp} over $k$ gives
$\sum_k\partial\Sv/\partial\vp{k} = \Mmat^{-1}(K_1+K_2+K_3) = \Mmat^{-1}\Mmat
= \mathbf{I}$.  Translating all three vertices by a common vector translates
$\Sv$ identically.  The implementation asserts this identity as a self-check.
\end{remark}

The implementation never forms $\Mmat^{-1}$ explicitly: it solves
$\Mmat\,X = K_k$ with \code{np.linalg.solve}, which is better conditioned near
the degenerate case and is natively batched over all triple points.

\begin{tcolorbox}[colback=blue!4!white, colframe=blue!40!black,
                  title={Code: first derivative of S}, fontupper=\small]
\codefile{src/partition/vectorized\_steiner.py}:
\code{\_compute\_dS\_dp(geom)} returns $\partial\Sv/\partial\vp{k}$;
\code{\_compute\_dS\_dlambda(geom, dS\_dp)} lifts it to $\lam$-space via
$\partial\Sv/\partial\lam_k = (\partial\Sv/\partial\vp{k})\,\ed{k}$.
\end{tcolorbox}

% =========================================================================
\section{First Derivatives of Perimeter and Area}
\label{sec:first}
% =========================================================================

\subsection{Steiner perimeter gradient}

Per cell per triple point the Steiner perimeter contribution is
$\ell_a + \ell_b - \ell_{ab}$, with
\begin{equation}
  \ell_a = \norm{\vp{a}-\Sv},\quad
  \ell_b = \norm{\vp{b}-\Sv},\quad
  \ell_{ab} = \norm{\vp{a}-\vp{b}},
\end{equation}
where $\vp{a},\vp{b}$ are the two variable points the cell borders.  Writing
$\mathbf{n}_a = (\vp{a}-\Sv)/\ell_a$ and using
$\partial\vp{a}/\partial\lam_i = \delta_{ai}\ed{i}$ and
$\partial\Sv/\partial\lam_i = (\Mmat^{-1}K_i)\,\ed{i}$,
\begin{equation}
  \pd{\ell_a}{\lam_i}
    = \mathbf{n}_a^\top
      \bigl(\delta_{ai}\,\ed{i} - (\Mmat^{-1}K_i)\,\ed{i}\bigr).
  \label{eq:dla}
\end{equation}
The chord term $\ell_{ab}$ does not involve $\Sv$; its gradient is the regular
perimeter-segment gradient of \codefile{docs/math/01-phase2-derivatives}.

\subsection{Steiner area Jacobian}

The Steiner area of a cell at a triple point is
$\Ast = \tfrac12\norm{(\vp{b}-\vp{a})\times(\Sv-\vp{a})} + \text{corner}$.
The corner triangle has a fixed mesh vertex and reuses the regular
$1$-inside area-Jacobian formula.  The \emph{void} triangle has all three
vertices moving, so write $\mathbf{e}_1=\vp{b}-\vp{a}$,
$\mathbf{e}_2=\Sv-\vp{a}$, $\mathbf{C}=\mathbf{e}_1\times\mathbf{e}_2$,
$\nhat=\mathbf{C}/\norm{\mathbf{C}}$.  Then
\begin{equation}
  \pd{\Ast_{\text{void}}}{\lam_i}
    = \tfrac12\,\nhat^\top
      \Bigl(\pd{\mathbf{e}_1}{\lam_i}\times\mathbf{e}_2
            + \mathbf{e}_1\times\pd{\mathbf{e}_2}{\lam_i}\Bigr),
  \label{eq:dA-void}
\end{equation}
the generalised scalar-triple-product derivative; the per-$\lam$ velocities
$\partial\mathbf{e}_1/\partial\lam_i$, $\partial\mathbf{e}_2/\partial\lam_i$
are built from $\ed{i}$ and $\partial\Sv/\partial\lam_i$.

\begin{tcolorbox}[colback=blue!4!white, colframe=blue!40!black,
                  title={Code: Steiner first derivatives}, fontupper=\small]
\codefile{src/partition/vectorized\_steiner.py}:
\code{compute\_steiner\_perimeter\_gradient\_analytical(pa)},
\code{compute\_steiner\_area\_jacobian\_analytical(pa)} and its \code{\_sparse}
twin.
\end{tcolorbox}

% =========================================================================
\section{Second Derivative of the Steiner Point}
\label{sec:second-S}
% =========================================================================

Differentiate \eqref{eq:dS-dp} with respect to $\vp{l}$.  Using
$\partial(\Mmat^{-1})/\partial\vp{l}
= -\Mmat^{-1}(\partial\Mmat/\partial\vp{l})\,\Mmat^{-1}$,
\begin{equation}
  \pdd{\Sv}{\vp{k}}{\vp{l}}
   = -\,\Mmat^{-1}\pd{\Mmat}{\vp{l}}\,\Mmat^{-1}K_k
     \;+\; \Mmat^{-1}\pd{K_k}{\vp{l}}.
  \label{eq:d2S}
\end{equation}
The two remaining ingredients are $\partial K_j/\partial\vp{l}$ and
$\partial\Mmat/\partial\vp{l} = \sum_j\partial K_j/\partial\vp{l}$.

Define the rank-two tensor
\begin{equation}
  T_{jl} \coloneqq \pd{\mathbf{u}_j}{\vp{l}}
        = \delta_{jl}\mathbf{I} - \pd{\Sv}{\vp{l}},
\end{equation}
which carries both the explicit dependence ($\vp{j}$ when $l=j$) and the
implicit dependence (through $\Sv$, always).  By Proposition~\ref{prop:unit},
\begin{equation}
  \pd{r_j}{\vp{l}} = \mathbf{n}_j^\top T_{jl},
  \qquad
  \pd{\mathbf{n}_j}{\vp{l}} = K_j\,T_{jl}.
\end{equation}
Differentiating $K_j = \Proj{\mathbf{n}_j}/r_j$ by the quotient rule,
\begin{equation}
  \pd{K_j}{\vp{l}}
   = -\frac{1}{r_j}\Bigl[
        (K_j T_{jl})\,\mathbf{n}_j^\top
        + \mathbf{n}_j\,(K_j T_{jl})^\top
        + K_j\,\pd{r_j}{\vp{l}}
      \Bigr].
  \label{eq:dK}
\end{equation}

\begin{remark}[Tensor structure of \eqref{eq:dK}]
\label{rem:rank3}
$\partial K_j/\partial\vp{l}$ is a rank-three object: two indices from the
$3\times3$ matrix $K_j$ and one from the differentiation direction $\vp{l}$.
The three terms of \eqref{eq:dK} are \emph{outer products} over the matrix and
derivative slots, \emph{not} matrix products --- $(K_j T_{jl})\,\mathbf{n}_j^\top$
means the entry $[a,b,c] = (K_j T_{jl})_{ac}\,(\mathbf{n}_j)_b$.  The
implementation builds each term with an explicit \code{np.einsum} index
string; a matrix product would silently contract the wrong axis.
\end{remark}

Substituting \eqref{eq:dK} (summed over $j$ for $\partial\Mmat/\partial\vp{l}$)
into \eqref{eq:d2S} gives $\partial^2\Sv/\partial\vp{k}\partial\vp{l}$, a
rank-three tensor per index pair $(k,l)$.

\begin{remark}[Second-order translation invariance]
\label{rem:translation2}
Differentiating Remark~\ref{rem:translation}'s identity
$\sum_k\partial\Sv/\partial\vp{k} = \mathbf{I}$ with respect to $\vp{l}$ gives
$\sum_k\partial^2\Sv/\partial\vp{k}\partial\vp{l} = \mathbf{0}$ for every $l$.
The implementation asserts this, together with the mixed-partial symmetry
$\partial^2\Sv/\partial\vp{k}\partial\vp{l}
= \partial^2\Sv/\partial\vp{l}\partial\vp{k}$ (with the corresponding swap of
the two input-component axes), as self-checks.
\end{remark}

% =========================================================================
\section{Chain Rule to \texorpdfstring{$\lam$}{lambda}-Space}
\label{sec:lambda}
% =========================================================================

Because each $\vp{k}$ is affine in its own $\lam_k$ with constant velocity
$\ed{k}$ and zero acceleration, the chain rule has no second-order vertex
term:
\begin{equation}
  \pdd{\Sv}{\lam_i}{\lam_j}
    = \Bigl(\pdd{\Sv}{\vp{i}}{\vp{j}}\Bigr)\cdot\ed{i}\cdot\ed{j},
  \label{eq:d2S-lambda}
\end{equation}
contracting the two input-component axes of the rank-three tensor with the
edge directions.

\begin{tcolorbox}[colback=blue!4!white, colframe=blue!40!black,
                  title={Code: second derivative of S}, fontupper=\small]
\codefile{src/partition/vectorized\_steiner.py}:
\code{\_compute\_d2S\_dp2(geom, dS\_dp)} builds
$\partial^2\Sv/\partial\vp{k}\partial\vp{l}$ from \eqref{eq:d2S}--\eqref{eq:dK};
\code{\_compute\_d2S\_dlambda(geom, d2S\_dp2)} applies \eqref{eq:d2S-lambda}.
\end{tcolorbox}

% =========================================================================
\section{Steiner Hessians}
\label{sec:hessians}
% =========================================================================

\subsection{Perimeter Hessian}

For a leg $\ell_a = \norm{\vp{a}-\Sv}$, write
$\boldsymbol{\Delta}_i = \partial(\vp{a}-\Sv)/\partial\lam_i
= \delta_{ai}\ed{i} - \partial\Sv/\partial\lam_i$ and
$\boldsymbol{\Delta}_{ij} = -\partial^2\Sv/\partial\lam_i\partial\lam_j$
(the variable point contributes no second-order term).  Then
\begin{equation}
  \pdd{\ell_a}{\lam_i}{\lam_j}
    = \boldsymbol{\Delta}_i^\top
      \frac{\Proj{\mathbf{n}_a}}{\ell_a}\,
      \boldsymbol{\Delta}_j
      \;+\; \mathbf{n}_a^\top\boldsymbol{\Delta}_{ij}.
  \label{eq:d2l}
\end{equation}
The first term is the projected dot product familiar from the regular
perimeter Hessian; the second term is the contribution of the curved
trajectory of $\Sv$ and is the reason the second derivative
\eqref{eq:d2S-lambda} is needed --- omitting it makes the Hessian wrong, not
merely approximate.  The chord $\ell_{ab}$ reuses the regular
perimeter-segment Hessian.  Each leg produces a full symmetric $3\times3$
block over the triple point's three variable points.

\subsection{Area Hessian}

For the void triangle area $\tfrac12\norm{\mathbf{C}}$,
$\mathbf{C}=\mathbf{e}_1\times\mathbf{e}_2$, $\nhat=\mathbf{C}/\norm{\mathbf{C}}$,
\begin{equation}
  \pdd{\Ast_{\text{void}}}{\lam_i}{\lam_j}
    = \tfrac12\Bigl[
        \Bigl(\pd{\mathbf{C}}{\lam_i}\Bigr)^{\!\top}
        \frac{\Proj{\nhat}}{\norm{\mathbf{C}}}
        \pd{\mathbf{C}}{\lam_j}
        + \nhat^\top\pdd{\mathbf{C}}{\lam_i}{\lam_j}
      \Bigr].
  \label{eq:d2A}
\end{equation}
The second-order cross-product term carries $\partial^2\mathbf{e}_2/
\partial\lam_i\partial\lam_j = \partial^2\Sv/\partial\lam_i\partial\lam_j$ ---
the area analogue of the perimeter's $\mathbf{n}_a^\top\boldsymbol{\Delta}_{ij}$
term.  The corner triangle has no moving mesh vertex, so its
$\partial^2\mathbf{C}$ retains only the mixed first-derivative cross terms.

\begin{tcolorbox}[colback=blue!4!white, colframe=blue!40!black,
                  title={Code: Steiner Hessians}, fontupper=\small]
\codefile{src/partition/vectorized\_steiner.py}:
\code{compute\_steiner\_perimeter\_hessian\_analytical(pa)},
\code{compute\_steiner\_area\_hessian\_analytical(pa, multipliers)}.  The
symmetric $3\times3$ blocks are scattered with \code{np.add.at} on the
pre-computed offset array \code{pa.tp\_hess\_off}.
\end{tcolorbox}

% =========================================================================
\section{Degenerate Case}
\label{sec:degenerate}
% =========================================================================

When some triangle angle reaches or exceeds $120^\circ$, the minimiser of
$G$ collapses onto the obtuse vertex: $\Sv = \vp{\text{obtuse}}$ identically.
The optimality condition \eqref{eq:fermat} no longer holds (three unit
vectors cannot make $120^\circ$ angles), and $\Mmat$ is singular there ---
$r_{\text{obtuse}}\to0$, so $K_{\text{obtuse}}$ diverges.  On the degenerate
branch the derivatives are elementary:
\begin{equation}
  \pd{\Sv}{\vp{k}} = \delta_{k,\text{obtuse}}\,\mathbf{I},
  \qquad
  \pdd{\Sv}{\vp{k}}{\vp{l}} = \mathbf{0}.
\end{equation}

\begin{remark}[Non-smoothness]
The two branches do not agree at the $120^\circ$ boundary --- $\Sv$ as a
function of $\Lam$ is genuinely non-smooth there.  The boundary is a
measure-zero set in parameter space, so the optimizer essentially never sits
on it, but each evaluation must commit to one branch.  The detection
(\code{degen\_mask}, \code{obtuse\_idx}) lives in the shared core
\code{\_compute\_tp\_steiner}, so the forward value $\Sv$ and its derivatives
are always evaluated on the same branch.  The non-degenerate formulas mask
out degenerate rows before inverting $\Mmat$.
\end{remark}

% =========================================================================
\section*{Summary table}
% =========================================================================

\begin{center}
\renewcommand{\arraystretch}{1.35}
\begin{tabular}{p{4.6cm}p{4.4cm}p{5.0cm}}
\toprule
\textbf{Quantity} & \textbf{Method} & \textbf{Implementing function} \\
\midrule
$\partial\Sv/\partial\vp{k}$
  & Analytical \eqref{eq:dS-dp}
  & \code{\_compute\_dS\_dp} \\
$\partial^2\Sv/\partial\vp{k}\partial\vp{l}$
  & Analytical \eqref{eq:d2S}--\eqref{eq:dK}
  & \code{\_compute\_d2S\_dp2} \\
$\partial\Sv/\partial\lam_i$,\;
$\partial^2\Sv/\partial\lam_i\partial\lam_j$
  & Chain rule \eqref{eq:d2S-lambda}
  & \code{\_compute\_dS\_dlambda}, \code{\_compute\_d2S\_dlambda} \\
Steiner perimeter gradient
  & Analytical \eqref{eq:dla}
  & \code{compute\_steiner\_perimeter\_gradient\_analytical} \\
Steiner area Jacobian
  & Analytical \eqref{eq:dA-void}
  & \code{compute\_steiner\_area\_jacobian\_analytical} \\
Steiner perimeter Hessian
  & Analytical \eqref{eq:d2l}
  & \code{compute\_steiner\_perimeter\_hessian\_analytical} \\
Steiner area Hessian
  & Analytical \eqref{eq:d2A}
  & \code{compute\_steiner\_area\_hessian\_analytical} \\
\bottomrule
\end{tabular}
\end{center}

\medskip
\noindent
The closed forms are validated against finite differences by
\codefile{testing/test\_steiner\_gradient\_analytical.py},
\codefile{testing/test\_steiner\_hessian\_analytical.py}, and
\codefile{testing/test\_steiner\_degenerate\_case.py}, and the assembled
Lagrangian Hessian by \codefile{testing/test\_exact\_hessian\_vs\_fd.py}.

\bibliographystyle{plain}
\bibliography{../shared/references}

\end{document}
